% GENERATED FILE. Edit canonical lessons, then run npm run generate:books.
% canonical-source-hash: fnv1a64:a068e32e41dce631
% canonical-lessons: PA-C118-macis, PA-C118-galica, PA-C118-sirhana, PA-C118-guddi, PA-C118-khidauna

\chapter{Matches, Carpet, Pillow, Doll, Toy}
\label{ch:pa-matches-carpet-pillow-doll-toy}
\hlchaptermodality{118}

\begin{chapteropening}
\textbf{By the end of this chapter:} \emph{I can say matches, a carpet, a pillow, a doll and a toy.}

Complete the last lesson of chapter 118: I can say matches, a carpet, a pillow, a doll and a toy.
\end{chapteropening}

\section[\texorpdfstring{\pa{ਮਾਚਿਸ} (mācis)}{ਮਾਚਿਸ (mācis)}]{\pa{ਮਾਚਿਸ} (mācis) — matches}

\label{lesson:PA-C118-macis}



\begin{quote}
Before the new one: say the Punjabi for a clay water pot, then the Punjabi for a candle.
\end{quote}

\subsection*{You'll want to know: \pa{ਮਾਚਿਸ}}
\textbf{\pa{ਮਾਚਿਸ}} — \emph{mācis} — \textquotedblleft{}matches\textquotedblright{}.

\begin{grammarlens}[title={What you've built}]
One more word for this chapter.
\end{grammarlens}

\subsection*{Your turn}
\begin{itemize}
\raggedright
  \item \emph{Say it:} \emph{mācis}
  \item \emph{Say it:} \emph{mācis}, once more
  \item \emph{Say it:} say \emph{mācis}
  \item \emph{Recall:} say the Punjabi for a clay water pot, then the Punjabi for a candle, then say \emph{mācis} again
\end{itemize}

\begin{tcolorbox}[breakable,colback=teal!4,colframe=teal!35!black,title={Before you move on}]
What does \pa{ਮਾਚਿਸ} mean? (\textquotedblleft{}Matches\textquotedblright{}.) Say it once more.
\end{tcolorbox}
\section[\texorpdfstring{\pa{ਗਲੀਚਾ} (galīcā)}{ਗਲੀਚਾ (galīcā)}]{\pa{ਗਲੀਚਾ} (galīcā) — a carpet}

\label{lesson:PA-C118-galica}



\begin{quote}
Before the new one: say the Punjabi for a candle, then the Punjabi for matches.
\end{quote}

\subsection*{You'll want to know: \pa{ਗਲੀਚਾ}}
\textbf{\pa{ਗਲੀਚਾ}} — \emph{galīcā} — \textquotedblleft{}a carpet\textquotedblright{}.

\begin{grammarlens}[title={What you've built}]
One more word for this chapter.
\end{grammarlens}

\subsection*{Your turn}
\begin{itemize}
\raggedright
  \item \emph{Say it:} \emph{galīcā}
  \item \emph{Say it:} \emph{galīcā}, once more
  \item \emph{Say it:} say \emph{galīcā}
  \item \emph{Recall:} say the Punjabi for a candle, then the Punjabi for matches, then say \emph{galīcā} again
\end{itemize}

\begin{tcolorbox}[breakable,colback=teal!4,colframe=teal!35!black,title={Before you move on}]
What does \pa{ਗਲੀਚਾ} mean? (\textquotedblleft{}A carpet\textquotedblright{}.) Say it once more.
\end{tcolorbox}
\section[\texorpdfstring{\pa{ਸਿਰਹਾਣਾ} (sirhāṇā)}{ਸਿਰਹਾਣਾ (sirhāṇā)}]{\pa{ਸਿਰਹਾਣਾ} (sirhāṇā) — a pillow}

\label{lesson:PA-C118-sirhana}



\begin{quote}
Before the new one: say the Punjabi for matches, then the Punjabi for a carpet.
\end{quote}

\subsection*{You'll want to know: \pa{ਸਿਰਹਾਣਾ}}
\textbf{\pa{ਸਿਰਹਾਣਾ}} — \emph{sirhāṇā} — \textquotedblleft{}a pillow\textquotedblright{}.

\begin{grammarlens}[title={What you've built}]
One more word for this chapter.
\end{grammarlens}

\subsection*{Your turn}
\begin{itemize}
\raggedright
  \item \emph{Say it:} \emph{sirhāṇā}
  \item \emph{Say it:} \emph{sirhāṇā}, once more
  \item \emph{Say it:} say \emph{sirhāṇā}
  \item \emph{Recall:} say the Punjabi for matches, then the Punjabi for a carpet, then say \emph{sirhāṇā} again
\end{itemize}

\begin{tcolorbox}[breakable,colback=teal!4,colframe=teal!35!black,title={Before you move on}]
What does \pa{ਸਿਰਹਾਣਾ} mean? (\textquotedblleft{}A pillow\textquotedblright{}.) Say it once more.
\end{tcolorbox}
\section[\texorpdfstring{\pa{ਗੁੱਡੀ} (guḍḍī)}{ਗੁੱਡੀ (guḍḍī)}]{\pa{ਗੁੱਡੀ} (guḍḍī) — a doll}

\label{lesson:PA-C118-guddi}



\begin{quote}
Before the new one: say the Punjabi for a carpet, then the Punjabi for a pillow.

Say \textquotedblleft{}fifth.\textquotedblright{} (\emph{Panjvāṁ}: \emph{panj} plus \emph{-vāṁ}, the ending every later ordinal takes.)
\end{quote}

\subsection*{You'll want to know: \pa{ਗੁੱਡੀ}}
\textbf{\pa{ਗੁੱਡੀ}} — \emph{guḍḍī} — \textquotedblleft{}a doll\textquotedblright{}.

\begin{grammarlens}[title={What you've built}]
One more word for this chapter.
\end{grammarlens}

\subsection*{Your turn}
\begin{itemize}
\raggedright
  \item \emph{Say it:} \emph{guḍḍī}
  \item \emph{Say it:} \emph{guḍḍī}, once more
  \item \emph{Say it:} say \emph{guḍḍī}
  \item \emph{Recall:} say the Punjabi for a carpet, then the Punjabi for a pillow, then say \emph{guḍḍī} again
\end{itemize}

\begin{tcolorbox}[breakable,colback=teal!4,colframe=teal!35!black,title={Before you move on}]
What does \pa{ਗੁੱਡੀ} mean? (\textquotedblleft{}A doll\textquotedblright{}.) Say it once more.
\end{tcolorbox}
\section[\texorpdfstring{\pa{ਖਿਡੌਣਾ} (khiḍauṇā)}{ਖਿਡੌਣਾ (khiḍauṇā)}]{\pa{ਖਿਡੌਣਾ} (khiḍauṇā) — a toy}

\label{lesson:PA-C118-khidauna}



\begin{quote}
Before the new one: say the Punjabi for a pillow, then the Punjabi for a doll.
\end{quote}

\subsection*{You'll want to know: \pa{ਖਿਡੌਣਾ}}
\textbf{\pa{ਖਿਡੌਣਾ}} — \emph{khiḍauṇā} — \textquotedblleft{}a toy\textquotedblright{}.

\begin{grammarlens}[title={What you've built}]
One more word for this chapter.
\end{grammarlens}

\subsection*{Your turn}
\begin{itemize}
\raggedright
  \item \emph{Say it:} \emph{khiḍauṇā}
  \item \emph{Say it:} \emph{khiḍauṇā}, once more
  \item \emph{Say it:} say \emph{khiḍauṇā}
  \item \emph{Recall:} say the Punjabi for a pillow, then the Punjabi for a doll, then say \emph{khiḍauṇā} again
\end{itemize}

\begin{tcolorbox}[breakable,colback=teal!4,colframe=teal!35!black,title={Before you move on}]
What does \pa{ਖਿਡੌਣਾ} mean? (\textquotedblleft{}A toy\textquotedblright{}.) Say it once more.
\end{tcolorbox}
